\newcommand{\cube}[4]{%
    \draw (#1,#2,#3) -- ++(#4,0,0) -- ++(0,#4,0) -- ++(-#4,0,0) -- cycle;
    \draw (#1,#2,#3) -- ++(0,0,#4);
    \draw (#1+#4,#2,#3) -- ++(0,0,#4);
    \draw (#1,#2+#4,#3) -- ++(0,0,#4);
    \draw (#1+#4,#2+#4,#3) -- ++(0,0,#4);
    \draw (#1,#2,#3+#4) -- ++(#4,0,0) -- ++(0,#4,0) -- ++(-#4,0,0) -- cycle;
}

\newcommand{\engl}[1]{(\textit{англ.} \textbf{#1})}

\newcommand{\cubesurface}[5]{%
    \fill[#5] (#1,#2,#3) -- ++(#4,0,0) -- ++(0,0,#4) -- ++(-#4,0,0) -- cycle;
    \draw (#1,#2,#3) -- ++(#4,0,0) -- ++(0,0,#4) -- ++(-#4,0,0) -- cycle;
    \fill[#5] (#1+#4,#2,#3) -- ++(0,#4,0) -- ++(0,0,#4) -- ++(0,-#4,0) -- cycle;
    \draw (#1+#4,#2,#3) -- ++(0,#4,0) -- ++(0,0,#4) -- ++(0,-#4,0) -- cycle;
    \fill[#5] (#1,#2,#3+#4) -- ++(#4,0,0) -- ++(0,#4,0) -- ++(-#4,0,0) -- cycle;
    \draw (#1,#2,#3+#4) -- ++(#4,0,0) -- ++(0,#4,0) -- ++(-#4,0,0) -- cycle;
}